\documentclass[11pt]{article}
\usepackage[T1]{fontenc}
\usepackage{lmodern}
\usepackage[utf8]{inputenc}
\usepackage[margin=1in]{geometry}
\usepackage{amsmath,amssymb,amsthm,mathtools}
\usepackage[hidelinks]{hyperref}

\newtheorem{definition}{Definition}
\newtheorem{theorem}{Theorem}
\newtheorem{proposition}{Proposition}
\newtheorem{remark}{Remark}
\newcommand{\Qsrc}{\mathcal Q_{\mathrm{src}}}
\newcommand{\Rel}{\mathsf R}
\newcommand{\Char}{\operatorname{Char}}
\newcommand{\supp}{\operatorname{supp}}

\title{Metric-Free Source Influence and Characteristic Target\\
\large P6-T01 registered draft/control formalization}
\author{The AEther-Flow Research Project}
\date{2026-07-26}

\begin{document}
\maketitle

\begin{abstract}
This registered draft/control manuscript executes v21 work item P6-T01.  It
defines a metric-free target for source-local influence, characteristic data,
and later causal qualification; proves a response-support composition
theorem; and gives an exact two-site countermodel family showing that the
unchanged P5 scalar flow and its diagonal tangent response do not select a
unique region-resolved influence relation.  The target is precise, but the
current ontology does not derive the localization, local evolution,
characteristic-extraction, orientation, operational, universality, or
robustness laws needed to realize it physically.  The result is therefore
\texttt{blocked\_adoption\_open\_continuation}, with P6-T02 permitted to
construct or obstruct one source-derived propagation structure.  No target
metric, GR null cone, physical causal cone, effective metric, source-law
adoption, or downstream promotion is supplied.
\end{abstract}

\section{Authority and source boundary}

The canonical ontology sources remain
\texttt{ontology/tex/aether\_flow\_foundations.tex},
\texttt{aether\_flow\_dynamics.tex}, and
\texttt{aether\_flow\_geometry.tex}.  The geometry manuscript is explicitly a
target-side exact-GR dictionary and states that no source-to-congruence or
source-to-metric bridge is supplied.  Its metric, null cones, proper time,
observer fields, and benchmark equations are therefore prohibited as premises
of the present source-side formalization.

P5-T08 supplies only the proposal-only one-amplitude family
\(\mathcal Q_{\mathrm{amp}}\), its conditional cubic forward flow
\[
 \phi_\tau(x)=\frac{x}{\sqrt{1+2\gamma\tau x^2}},
 \qquad \gamma>0,\quad \tau\geq0,
\]
the exact diagonal tangent response
\[
 D\phi_\tau(x)=(1+2\gamma\tau x^2)^{-3/2},
\]
and scoped obstruction data.  The parameter \(\tau\) is not identified with
physical time.  The scoped historical
\(g_{\mathrm{eff}}^{\mathrm{GSC\mbox{-}cand}}\) record remains exactly within
its declared source-extension status and is not a physical Lorentzian metric
or a premise here.

\section{Metric-free target}

\begin{definition}[Source-local response datum]
\label{def:datum}
A \emph{source-local response datum} is a tuple
\[
 \mathfrak C_{\mathrm{src}}=
 (\mathfrak B,\Qsrc,\operatorname{res},\operatorname{glue},
 \Phi_{\mathrm{src}},D\Phi_{\mathrm{src}},\{V_i(q)\}_{i\in I}),
\]
with the following roles.
\begin{enumerate}
  \item \(\mathfrak B\) is a source-defined localization base or finite region
  index set \(I\).  It is not a target spacetime atlas.
  \item \(\Qsrc(U)\) gives admissible source states on a source region \(U\),
  with declared restriction and gluing operations.
  \item \(\Phi_{\mathrm{src},\tau}\) is a declared local source evolution and
  \(D\Phi_{\mathrm{src},\tau}(q)\) its linearized response.  The parameter and its
  composition law must be derived or explicitly proposed on the source side.
  \item At a finite test state \(q\), the tangent data decompose as
  \(T_q\Qsrc=\bigoplus_{i\in I}V_i(q)\).  This finite block form is a test
  interface; it is not asserted to exist in the current ontology.
\end{enumerate}
No metric, norm, angle, target coordinate, or GR causal structure occurs in
this definition.
\end{definition}

\begin{definition}[Formal influence relation]
\label{def:influence}
For a datum in Definition~\ref{def:datum}, let
\[
 L_{\tau,q}=D\Phi_{\mathrm{src},\tau}(q),\qquad
 L_{\tau,q}^{\,ji}:V_i(q)\longrightarrow
 V_j(\Phi_{\mathrm{src},\tau}q)
\]
be its block response.  Define
\[
 i\;\Rel_{\tau,q}\;j
 \quad\Longleftrightarrow\quad
 L_{\tau,q}^{\,ji}\neq0.
\]
Thus \(i\Rel_{\tau,q}j\) means that some infinitesimal source perturbation in
block \(i\) has nonzero linear response in block \(j\).  This is a
\emph{formal source influence} relation.  It is not yet physical causality.
\end{definition}

\begin{definition}[Infinitesimal and characteristic targets]
\label{def:characteristic}
If a source generator \(X\) exists, its infinitesimal influence relation is
defined from the nonzero blocks of \(DX(q)\).  If, in addition, a source-local
differential equation \(\mathcal E(q)=0\) of order \(m\) and an admissible
square reduced principal symbol \(\sigma_q(\xi)\) are source-derived, define
\[
 \Char_q^*=
 \{\xi\neq0:\det\sigma_q(\xi)=0\}.
\]
Equivalently, a source-derived scalar principal polynomial
\(P_q(\xi)\) may define \(\Char_q^*=\{P_q=0\}\).  If \(P_q\) is hyperbolic
with respect to a source-derived covector \(h\), its hyperbolicity component
\(\Gamma_h(P_q)\) and the appropriate dual propagation cone are eligible
\emph{characteristic candidates}.  Their orientation, branch selection,
nondegeneracy, covariance, robustness, and operational signal meaning remain
separate obligations.
\end{definition}

\begin{remark}[Causal qualification]
The word \emph{causal} is reserved for a formal influence or characteristic
candidate that also passes the Gate B obligations: source provenance,
orientation, local consistency, hyperbolicity where applicable,
nondegeneracy, naturality, robustness, operational signal meaning, and
single-structure universality across admissible signal sectors.  Metric scale
is intentionally deferred.  A GR null cone or target metric cannot satisfy
these obligations by being copied into source notation.
\end{remark}

\section{Response-support composition}

\begin{theorem}[P6T01-THM-RESPONSE-SUPPORT-COMPOSITION]
\label{thm:composition}
Assume a finite source-local response datum and the response cocycle
\[
 L_{\tau+\sigma,q}
 =L_{\sigma,\Phi_{\mathrm{src},\tau}q}\,L_{\tau,q}.
\]
Then
\[
 \Rel_{\tau+\sigma,q}
 \subseteq
 \Rel_{\sigma,\Phi_{\mathrm{src},\tau}q}\circ\Rel_{\tau,q}.
\]
In words, every nonzero block of a composed response has at least one
intermediate block through which both component responses are nonzero.
Equality need not hold because different intermediate contributions can
cancel.
\end{theorem}

\begin{proof}
The \(ki\) block of the cocycle is
\[
 L_{\tau+\sigma,q}^{\,ki}
 =\sum_{j\in I}
 L_{\sigma,\Phi_{\mathrm{src},\tau}q}^{\,kj}
 L_{\tau,q}^{\,ji}.
\]
If the left side is nonzero, at least one summand is nonzero.  For that
intermediate \(j\), both factors are nonzero, so
\(i\Rel_{\tau,q}j\) and
 \(j\Rel_{\sigma,\Phi_{\mathrm{src},\tau}q}k\).  This proves the inclusion.  The reverse
inclusion can fail by cancellation and is not asserted.
\end{proof}

\section{Exact two-site nonselection countermodel}

\begin{definition}[Proposal-only two-site extension family]
\label{def:twosite}
For \(\gamma>0\) and \(c\geq0\), consider the source-side finite system
\[
\begin{aligned}
 \dot a&=-\gamma a^3+c(b-a),\\
 \dot b&=-\gamma b^3+c(a-b).
\end{aligned}
\tag{1}\label{eq:twosite}
\]
The two sites and the coupling \(c\) are proposal-only source-extension data
used as a finite countermodel family.  They are not derived from or adopted
by the current ontology.
\end{definition}

\begin{theorem}[P6T01-THM-DIAGONAL-RESPONSE-NONSELECTION]
\label{thm:nonselection}
For every \(c\geq0\), the diagonal \(a=b=x\) is invariant under
\eqref{eq:twosite}, and its restricted evolution is exactly
\(\dot x=-\gamma x^3\), the P5 cubic candidate.  Along the diagonal solution
\(x(\tau)=\phi_\tau(x_0)\), the two-site tangent response is
\[
 \mathcal R_c(\tau;x_0)
 =
 D\phi_\tau(x_0)\,
 \frac12
 \begin{pmatrix}
  1+e^{-2c\tau} & 1-e^{-2c\tau}\\
  1-e^{-2c\tau} & 1+e^{-2c\tau}
 \end{pmatrix}.
\tag{2}\label{eq:response}
\]
Consequently:
\begin{enumerate}
  \item the diagonal response on the symmetric perturbation \((1,1)\) is
  \(D\phi_\tau(x_0)(1,1)\) for every \(c\);
  \item when \(c=0\), both cross-site response blocks vanish;
  \item when \(c>0\) and \(\tau>0\), both cross-site response blocks are
  strictly positive.
\end{enumerate}
Thus the exact P5 scalar evolution and diagonal tangent response are
compatible with distinct region-resolved influence relations.  They do not
select a unique one.
\end{theorem}

\begin{proof}
Substituting \(a=b=x\) into \eqref{eq:twosite} cancels the coupling terms and
gives \(\dot x=-\gamma x^3\).  The Jacobian along the diagonal is
\[
 A_c(\tau)=
 -3\gamma x(\tau)^2 I+
 cK,\qquad
 K=
 \begin{pmatrix}-1&1\\1&-1\end{pmatrix}.
\]
The scalar multiple of \(I\) commutes with \(K\) at all parameter values.
Moreover,
\[
 \exp\!\left(-\int_0^\tau3\gamma x(s)^2\,ds\right)
 =D\phi_\tau(x_0),
\]
while \(K\) has eigenvalue \(0\) on \((1,1)\) and eigenvalue \(-2\) on
\((1,-1)\).  Therefore the fundamental response is
\(D\phi_\tau(x_0)e^{c\tau K}\), which is exactly
\eqref{eq:response}.  The three consequences follow immediately.
\end{proof}

\begin{proposition}[P6T01-PROP-PRECISE-SCOPED-OBSTRUCTION]
\label{prop:obstruction}
The current ontology and unchanged P5 input do not derive a unique
region-resolved influence relation, characteristic operator, principal
polynomial, hyperbolicity cone, time orientation, or operational causal
structure.  In particular, Theorem~\ref{thm:nonselection} blocks the exact
route that would infer a unique local influence relation solely from the P5
diagonal flow and tangent response.
\end{proposition}

\begin{proof}
The \(c=0\) and \(c>0\) members have the same restricted P5 flow and the same
diagonal response, but different cross-site influence relations.  Therefore
those P5 data do not select a unique local relation.  The remaining listed
objects require additional source definitions and laws not present in the
input.  This conclusion is scoped to the unchanged P5 data and the displayed
finite extensions; it does not rule out conservative source extensions or
prove a global no-go theorem.
\end{proof}

\section{Controls and Gate B}

The task-local Gate B checklist records four required controls.
\begin{enumerate}
  \item \textbf{Positive influence control:} \(c>0\) makes cross-site formal
  influence detectable, while failing source-provenance and physical-
  interpretation gates because the sites and \(c\) were proposed for the
  control.
  \item \textbf{Degenerate control:} \(c=0\) preserves the same diagonal P5
  response but has no cross-site influence.
  \item \textbf{Acausal control:} a synthetic relation that violates an
  asserted source-derived orientation must fail orientation consistency.  It
  is a test fixture, not a claim that \eqref{eq:twosite} is physically
  acausal.
  \item \textbf{Multi-cone control:} two inequivalent hyperbolicity branches
  with no source selector must fail uniqueness and common-signal
  universality.
\end{enumerate}

The metric-free target definition, response-support theorem, controls, and
no-target-import check pass at P6-T01.  Gate B as a physical reconstruction
gate does not pass: localization provenance, restriction and gluing, local
source law, characteristic extraction, orientation, hyperbolicity,
naturality, robustness, operational signal meaning, and universality are
absent or untested.  Scale remains a later burden by design.

\section{Disposition and next bounded route}

P6-T01 is complete as a target formalization with a precise scoped
nonselection obstruction.  Adoption remains blocked and same-milestone
continuation remains open:
\[
\texttt{blocked\_adoption\_open\_continuation}.
\]
This is not a global no-go theorem.
The exact branch ``infer a unique region influence relation from the
unchanged P5 diagonal response'' is locally frozen.  P6-T02 may instead make
one bounded Candidate Constructor attempt to supply a source-derived
propagation relation, principal polynomial, characteristic order, or a new
precise obstruction.  Such a packet must add source-side localization or law
content and cannot reopen the frozen inference by relabeling it.

\section{Explicit nonclaims}

This manuscript does not:
\begin{itemize}
  \item edit, adopt, reject, or modify canonical ontology or a source law;
  \item identify \(\tau\) as physical time or the two-site family as physical
  spacetime;
  \item construct or adopt a physical causal relation, characteristic cone,
  conformal class, \(M_{\mathrm{src}}\), or unscoped \(g_{\mathrm{eff}}\);
  \item alter the protected scoped status of the historical
  \(g_{\mathrm{eff}}^{\mathrm{GSC\mbox{-}cand}}\) record;
  \item import target coordinates, a GR null cone, a Lorentzian metric,
  proper time, detector semantics, matter coupling, or Einstein equations;
  \item prove a global no-go theorem, reject the program, or close future
  source extensions;
  \item promote the exact-GR benchmark or create proof, Gate Chair,
  publication, push, or completed-derivation authority.
\end{itemize}

\end{document}
